\documentclass{article}
\usepackage[margin=25mm]{geometry}
\usepackage{amsmath,amssymb,amsthm}
\usepackage{xurl}
\usepackage[hidelinks]{hyperref}
% Shared commands for notes in docs/paper-gaps/.

\DeclareUnicodeCharacter{2080}{\ifmmode{}_{0}\else\textsubscript{0}\fi}
\DeclareUnicodeCharacter{2081}{\ifmmode{}_{1}\else\textsubscript{1}\fi}
\DeclareUnicodeCharacter{2082}{\ifmmode{}_{2}\else\textsubscript{2}\fi}
\DeclareUnicodeCharacter{2099}{\ifmmode{}_{n}\else\textsubscript{n}\fi}

\newcommand{\Lean}{\textsc{Lean}}
\ExplSyntaxOn
\NewDocumentCommand{\leanid}{m}{
  \ifmmode
    \textsf{\detokenize{#1}}
  \else
    \group_begin:
    \urlstyle{sf}
    \tl_set:Nn \l_tmpa_tl {#1}
    \tl_replace_all:Nnn \l_tmpa_tl {\_} {_}
    \exp_args:NV \path \l_tmpa_tl
    \group_end:
  \fi
}
\ExplSyntaxOff
\providecommand{\tr}{\operatorname{tr}}
\newcommand{\tnleanIssueBase}{https://github.com/LionSR/TNLean/issues/}
\newcommand{\tnleanPrBase}{https://github.com/LionSR/TNLean/pull/}
\newcommand{\tnleanDocBase}{https://sirui-lu.com/TNLean/docs/}
\newcommand{\ghissue}[1]{\href{\tnleanIssueBase#1}{GitHub issue \##1}}
\newcommand{\ghpr}[1]{\href{\tnleanPrBase#1}{PR \##1}}
\newcommand{\doclink}[1]{\href{\tnleanDocBase#1.html}{\path{#1}}}

% Dirac notation and the identity operator, as in blueprint/src/macros/common.tex.
% \providecommand keeps note-local definitions authoritative.
\usepackage{dsfont}
\providecommand{\ket}[1]{|#1\rangle}
\providecommand{\bra}[1]{\langle #1|}
\providecommand{\braket}[2]{\langle #1 | #2 \rangle}
\providecommand{\ketbra}[2]{|#1\rangle\!\langle #2|}
\providecommand{\Id}{\mathds{1}}
\providecommand{\one}{\mathds{1}}
\providecommand{\C}{\mathbb{C}}
\providecommand{\MN}[1]{M_{#1}(\C)}
\providecommand{\MD}{M_D(\C)}

% Verdict marker: \gapnote{<kind>}{<status>}. Kind is the policy
% classification (clarification, local-correction, scope-restriction,
% unfaithful, false-source, open-gap); status is open, wip (elimination
% underway), resolved, or historical. Machine-read by the site index;
% prints a verdict line.
\newcommand{\gapnote}[2]{\par\noindent\textbf{Verdict:} #1 (#2).\par}

\providecommand{\gapnote}[2]{\par\noindent\textbf{Verdict:} #1 (#2).\par}
\newtheorem{proposition}{Proposition}
\providecommand{\C}{\mathbb C}
\providecommand{\cH}{\mathcal H}
\newcommand{\vac}{\mathrm{vac}}
\newcommand{\dep}{\mathrm{depth}}
\title{A Tree Counterexample to the Unrestricted Summed Basis Bound in Hastings}
\author{}
\date{October 28, 2026}
\begin{document}
\maketitle
\gapnote{false-source}{resolved}

\section{The assertion and the counterexample}

Hastings, \emph{Solving gapped Hamiltonians locally}~\cite{Hastings2006Solving},
\href{https://arxiv.org/abs/cond-mat/0508554v3}{arXiv:cond-mat/0508554v3},
source \path{proj3new.tex}, lines 169--190, assumes a metric lattice,
interaction range $R$, term norm bound $J$, and at most $S$ sites within
distance $R$ of a site. Appendix B assumes a ground-space dimension and a
positive gap bounded uniformly in system size (lines 1535--1538).
For the localized terms $M_i$, lines 1678--1694 assert that one can choose a
ground-space basis with $\sum_i o(i)$ bounded independently of system size,
where $o(i)$ is the norm of the off-diagonal part of $PM_iP$. The printed
bound (oibu) also contains the reciprocal of the largest compressed
spectral spread.

The following construction disproves this assertion for general metric
lattices. It does not disprove the version with polynomial volume growth.

\begin{proposition}
For every integer $L\ge2$ there is a qubit Hamiltonian
$H_L=\sum_{v\in T_L}h_v$ on a finite rooted tree, with its graph metric, such
that:
\begin{enumerate}
\item every $h_v$ is Hermitian, supported in the radius-one ball about $v$,
  and has norm at most $769$; every radius-one ball has at most $258$ sites;
\item the ground energy is zero, the ground space has dimension two, and
  the spectral gap is at least $735/2048$;
\item for the actual normalized Gaussian construction (mid) with
  $l_{\rm proj}=R=1$, one has $M_v=h_v$ for every positive filter parameter;
\item in every orthonormal ground-space basis,
  \begin{equation}
    \sum_{v\in T_L}o(v)\ge2^{L-2};
    \label{eq:hastings05_tree_basis_counterexample_lower}
  \end{equation}
  the largest compressed spectral spread belongs to the root and lies
  between $3$ and $4$.
\end{enumerate}
The number of sites is comparable to $256^L$, so the lower bound is of
order $N_L^{1/8}$.
\end{proposition}

\section{The positive Hamiltonian and its one-particle operator}

Put $b=256$ and $a=1/128$. Let $T_L$ be the complete rooted $b$-ary tree
of depth $L$, with root $r$. Thus
\[
 N_L=\sum_{j=0}^Lb^j=\frac{b^{L+1}-1}{b-1}.
\]
On $\cH_L=\bigotimes_{v\in T_L}\C^2$, write
$s_v=|0\rangle\langle1|_v$ and $n_v=s_v^*s_v$.
The occupation vectors $|A\rangle$, indexed by subsets $A\subseteq T_L$,
form an orthonormal basis. The number operator
$\widehat N=\sum_v n_v$ has eigenspaces $\cH_m$ spanned by the subsets of
cardinality $m$.
For an oriented edge $(p,c)$ from parent to child, set
\begin{equation}
 d_{pc}=s_c-a s_p,\qquad
 H_L=\sum_{(p,c)}d_{pc}^*d_{pc}.
 \label{eq:hastings05_tree_basis_counterexample_H}
\end{equation}
This Hamiltonian is positive and preserves every number sector.

Identify the one-particle space with $\ell^2(T_L)$, and define
\[
 D:\ell^2(T_L)\longrightarrow\ell^2(T_L\setminus\{r\}),\qquad
 (Df)(c)=f(c)-a f(p(c)).
\]
The one-particle restriction of $H_L$ is $D^*D$.
Write $D^*=J-aP$, where $Jy$ places $y_c$ at each non-root vertex and zero
at the root, and $(Py)_p=\sum_{c:\,p(c)=p}y_c$.
The map $J$ is isometric. Cauchy--Schwarz, followed by the disjointness of
the sets of children, gives
\[
 \|Py\|^2\le b\sum_c|y_c|^2.
\]
Consequently
\begin{equation}
 \|D^*y\|\ge(1-a\sqrt b)\|y\|=\frac78\|y\|.
 \label{eq:hastings05_tree_basis_counterexample_Dgap}
\end{equation}
The kernel of $D$ is one-dimensional, since its equations require
$f(v)=a^{\dep(v)}f(r)$. Its positive normalized generator is
\[
 \psi_v=c_La^{\dep(v)},\qquad
 c_L^2=\left(\sum_{j=0}^L(ba^2)^j\right)^{-1}.
\]
Since $ba^2=1/64$,
\begin{equation}
 \frac{63}{64}\le c_L^2\le1,
 \qquad \sum_{v\ne r}\psi_v^2\le\frac1{64}.
 \label{eq:hastings05_tree_basis_counterexample_tail}
\end{equation}
By \eqref{eq:hastings05_tree_basis_counterexample_Dgap},
$DD^*\ge(49/64)I$. The nonzero eigenvalues of $D^*D$ and $DD^*$ coincide:
an eigenvector of either matrix with positive eigenvalue is carried to
one of the other by $D$ or $D^*$. Hence, with $\delta=49/64$,
\begin{equation}
 D^*D\ge\delta(I-|\psi\rangle\langle\psi|).
 \label{eq:hastings05_tree_basis_counterexample_oneparticle}
\end{equation}

\section{A gap in every number sector}

For a coefficient matrix $A$ on $\ell^2(T_L)$ define
\[
 \mathcal L(A)=\sum_{u,v}A_{uv}s_u^*s_v.
\]
This linear map preserves positivity. Indeed, if $A=C^*C$, then
\[
 \mathcal L(A)=\sum_k\left(\sum_v C_{kv}s_v\right)^*
                         \left(\sum_v C_{kv}s_v\right)\ge0.
\]
This argument uses only products in the indicated order; it makes no
assumption of canonical bosonic or fermionic commutation relations.
Since the coordinates of $\psi$ are real, putting $B=\sum_v\psi_vs_v$
in \eqref{eq:hastings05_tree_basis_counterexample_oneparticle} gives
\begin{equation}
 H_L\ge\delta(\widehat N-B^*B).
 \label{eq:hastings05_tree_basis_counterexample_lift}
\end{equation}

For $\Phi\in\cH_m$ one has
$\sum_v\|s_v\Phi\|^2=m\|\Phi\|^2$. Vector-valued Cauchy--Schwarz and
\eqref{eq:hastings05_tree_basis_counterexample_tail} therefore imply
\[
 \|B\Phi\|\le\|\psi_rs_r\Phi\|
       +\left\|\sum_{v\ne r}\psi_vs_v\Phi\right\|
       \le\left(1+\frac{\sqrt m}{8}\right)\|\Phi\|.
\]
For $m\ge2$, the elementary inequality $1/\sqrt m\le3/4$ gives
$\|B\Phi\|\le(7/8)\sqrt m\,\|\Phi\|$.
Combining this with \eqref{eq:hastings05_tree_basis_counterexample_lift},
\[
 \langle\Phi,H_L\Phi\rangle
 \ge\frac{49}{64}\frac{15m}{64}\|\Phi\|^2
 \ge\frac{735}{2048}\|\Phi\|^2\qquad(m\ge2).
\]
In the one-particle sector, the kernel is $\C\psi$ and its orthogonal
complement has energy at least $49/64$. The zero-particle sector is the
vacuum and has energy zero. The number sectors are orthogonal invariant
subspaces, so the ground space is exactly
\begin{equation}
 \mathcal G_L=\operatorname{span}\{|\vac\rangle,|\psi\rangle\},
 \label{eq:hastings05_tree_basis_counterexample_ground}
\end{equation}
and the asserted uniform many-body gap follows.

\section{A local decomposition with cancelling fields}

Let $X_v,Y_v$ denote the usual Pauli matrices, with
$Y=\left(\begin{smallmatrix}0&-i\\i&0\end{smallmatrix}\right)$.
Put $V_v=X_v$ at even depth and $V_v=Y_v$ at odd depth. Define
\[
 T_r=-\sum_{c:\,p(c)=r}V_c,\qquad
 T_v=V_v-\sum_{c:\,p(c)=v}V_c\quad(v\ne r),
\]
and
\begin{equation}
 h_v=\sum_{c:\,p(c)=v}d_{vc}^*d_{vc}+T_v.
 \label{eq:hastings05_tree_basis_counterexample_terms}
\end{equation}
Every non-root field occurs once with each sign, so $\sum_vT_v=0$ and
$\sum_vh_v=H_L$. Each $h_v$ is Hermitian and supported at $v$ and its
children, hence in the radius-one ball about $v$. There are at most
$b+2=258$ sites in any radius-one ball. Moreover,
$\|d_{vc}^*d_{vc}\|\le(1+a)^2<2$ and $\|T_v\|\le b+1$, whence
\[
 \|h_v\|\le2b+(b+1)=769.
\]
The source imposes no positivity condition on the individual interaction
terms. Thus these cancelling fields do not violate its hypotheses.

Let $P_0$ be the projection onto \eqref{eq:hastings05_tree_basis_counterexample_ground}.
In the ordered basis $(|\vac\rangle,|\psi\rangle)$, put
$\sigma_j=\sigma_x$ for even $j$ and $\sigma_j=\sigma_y$ for odd $j$.
The edge operators annihilate both ground vectors, and a direct matrix
element calculation gives
$P_0V_vP_0=c_La^{\dep(v)}\sigma_{\dep(v)}$.
The compressed interaction terms are therefore
\begin{align}
 P_0h_rP_0&=-2c_L\sigma_y,\label{eq:hastings05_tree_basis_counterexample_root}\\
 P_0h_vP_0&=c_La^j(\sigma_j-2\sigma_{j+1})
       &&(1\le j=\dep(v)<L),\label{eq:hastings05_tree_basis_counterexample_internal}\\
 P_0h_vP_0&=c_La^L\sigma_L
       &&(\dep(v)=L).\label{eq:hastings05_tree_basis_counterexample_leaf}
\end{align}
The spectral spread of the root compression is $4c_L\in[3,4]$.
For another internal vertex it is $2\sqrt5\,c_La^j<4c_L$, since
the two Pauli matrices anticommute; at a leaf it is $2c_La^L<4c_L$.
Thus the largest compressed spectral spread is uniformly positive and
is attained at the root.

\section{The lower bound in every ground-space basis}

For an orthonormal basis $\mathcal B$ of the two-dimensional ground space,
write $o_{\mathcal B}(A)$ for the operator norm of its off-diagonal part.
This operation is linear before taking the norm, so $o_{\mathcal B}$ is
subadditive and absolutely homogeneous. A qubit basis has a unit Bloch
axis $u\in\mathbb R^3$. Its off-diagonal norms satisfy
\[
 o_{\mathcal B}(\sigma_x)=\sqrt{1-u_x^2},\qquad
 o_{\mathcal B}(\sigma_y)=\sqrt{1-u_y^2},\qquad
 o_{\mathcal B}(\sigma_x)+o_{\mathcal B}(\sigma_y)\ge1.
\]
For the last inequality, the sum of the squares of the two nonnegative
terms is $2-u_x^2-u_y^2\ge1$.
Put $P=\sigma_L$, $Q=\sigma_{L-1}$ and
$x=o_{\mathcal B}(P)$, $y=o_{\mathcal B}(Q)$,
$z=o_{\mathcal B}(Q-2P)$. Then $x+y\ge1$ and $y\le z+2x$.
The quantity $A=2x+z$ is at least $2x$ and at least $y$, so
$3A\ge2(x+y)\ge2$.
There are $b^L$ leaves and $b^{L-1}$ vertices in the preceding level.
Since $L\ge2$, this preceding level does not include the root, and
\eqref{eq:hastings05_tree_basis_counterexample_internal} applies there.
Keeping only these two levels in the nonnegative sum gives
\begin{align*}
 \sum_v o_{\mathcal B}(P_0h_vP_0)
 &\ge c_L(ba)^Lx+c_L(ba)^{L-1}z\\
 &=c_L2^{L-1}(2x+z)
 \ge\frac23c_L2^{L-1}\ge2^{L-2},
\end{align*}
where the last step uses $c_L\ge3/4$.
Finally $b^L\le N_L\le2b^L$ and $b=2^8$, so this bound grows at least
as a positive constant times $N_L^{1/8}$.

\section{The actual localized operators and the source scope}

Take $R=l_{\rm proj}=1$. In the graph metric, the source's truncated
generator at $v$ is exactly
\[
 H_{{\rm loc},v}=\sum_{w:\,d(v,w)\le l_{\rm proj}-R}h_w=h_v.
\]
For every $q>0$ and every $\Delta E>0$, its normalized Gaussian filter is
\[
 M_v=\frac{\Delta E}{\sqrt{2\pi q}}
       \int_{\mathbb R}e^{-(t\Delta E)^2/(2q)}
           e^{ith_v}h_ve^{-ith_v}\,dt=h_v.
\]
Indeed the conjugated operator is constant and the scalar Gaussian has
integral one. This includes the source choice $q=c_1l_{\rm proj}\Delta E$
for any fixed $c_1>0$. The normalization is the intended one of (thd);
the missing $q$ in the printed (mid) is discussed in
\path{hastings05_truncated_dynamics.tex}.
Thus \eqref{eq:hastings05_tree_basis_counterexample_lower} concerns the
actual $M_v$, not a replacement by first filters under $H_L$.

All local bounds, the ground-space dimension, and a positive lower bound
on the gap are independent of $L$. A Lieb--Robinson time-window slope may
also be chosen uniformly from these local bounds. The maximal compressed
spectral spread is bounded away from zero. Nevertheless the summed
off-diagonal norm diverges in every basis.
Appendix B does not impose a lower bound on $l_{\rm proj}/l_{\rm fac}$
before asserting (oibu) at lines 1678--1694. The later choice of a
size-dependent projection radius for a good global approximation is a
different requirement. No counterexample along that coupled choice of
radius is asserted here.

\section{Formalization status and corrected statement}

The family and its many-body gap proof above are mathematical arguments;
they have not been formalized in Lean. The verdict records the failure of
the unrestricted printed claim, not a kernel-verified counterexample.
The separately verified construction of a basis under polynomial
sphere growth remains valid. That additional hypothesis makes the
exponential anchor sums finite uniformly in system size, as recorded in
\path{hastings05_geometric_growth_scope.tex} and
\path{hastings05_two_dim_ground_basis.tex}.
The exponential truncated-term estimate (facld) is a separate statement;
this counterexample does not refute it. The actual physical metric truncation
estimate is now proved without polynomial growth, by the weighted-tail
argument documented in \path{hastings05_geometric_growth_scope.tex}.

\bibliographystyle{unsrt}
\bibliography{references}
\end{document}
